\documentclass[a4paper]{article}
\usepackage{graphicx}
\usepackage{amsmath,amssymb}
\usepackage{hyperref}
\usepackage{minted}
\usepackage{multirow}
\usepackage{float}
\usepackage{todonotes}
\usepackage{forest}
\usepackage{xstring}
\input{/home/donkarlo/Dropbox/repo/nd_language_ai_project/src/nd_language_ai/formal/computer/programming/tex/latex/code_clips/external_dot.tex}
\title{Ermittlung des Teilergebnisses „Sprechen“}
\date{}

\begin{document}
    \maketitle
        
        \textbf{Deutsch:}
        Die Bewertung des Subtests „Sprechen“ erfolgt durch zwei lizenzierte Prüfer bzw. Prüferinnen gemäß den Bewertungskriterien. Bei den Stufen B1 und A2 besteht die Möglichkeit zu unterscheiden, ob die Kriterien „gut erfüllt“ oder „erfüllt“ wurden. Wichtig ist aber stets, dass die Prüfer bzw. Prüferinnen ihr Urteil kriterienbasiert und nicht nach Punktwerten fällen.
        
        \textbf{English:}
        The evaluation of the subtest “Speaking” is carried out by two licensed examiners according to the assessment criteria. At levels B1 and A2 it is possible to distinguish whether the criteria are “well fulfilled” or “fulfilled”. However, it is important that the examiners base their judgement on criteria rather than purely on point values.
        
        \begin{table}[H]
            \centering
            \begin{tabular}{|l|c|c|c|c|c|c|}
                \hline
                & \multicolumn{2}{c|}{B1} & \multicolumn{2}{c|}{A2} & A1 & 0 \\
                \cline{2-7}
                & gut erfüllt & erfüllt & gut erfüllt & erfüllt & erfüllt &   \\
                \hline
                \multicolumn{7}{|l|}{\textbf{I Aufgabenbewältigung (Task accomplishment)}} \\
                \hline
                Teil 1A                                    & 5           & 4       & 3           & 2       & 1       & 0 \\
                Teil 1B                                    & 5           & 4       & 3           & 2       & 1       & 0 \\
                Teil 2A                                    & 10          & 8       & 6           & 4       & 2       & 0 \\
                Teil 2B                                    & 10          & 8       & 6           & 4       & 2       & 0 \\
                Teil 3                                     & 20          & 16      & 12          & 8       & 4       & 0 \\
                \hline
                Summe 1                                    & 50          & 40      & 30          & 20      & 10      & 0 \\
                \hline
                II Aussprache / Intonation (Pronunciation) & 10          & 8       & 6           & 4       & 2       & 0 \\
                III Flüssigkeit (Fluency)                  & 10          & 8       & 6           & 4       & 2       & 0 \\
                IV Korrektheit (Correctness)               & 15          & 12      & 9           & 6       & 3       & 0 \\
                V Wortschatz (Vocabulary)                  & 15          & 12      & 9           & 6       & 3       & 0 \\
                \hline
                Summe 2                                    & 50          & 40      & 30          & 20      & 10      & 0 \\
                \hline
                Gesamtergebnis (Summe 1+2)                 & 100         & 80      & 60          & 40      & 20      & 0 \\
                \hline
            \end{tabular}
        \end{table}
        
        Für das Erreichen der Stufen A2 und B1 gilt:
        
        \textbf{English:} The following score ranges apply for reaching levels A2 and B1.
        
        \begin{table}[H]
            \centering
            \begin{tabular}{|c|c|}
                \hline
                Punkte (Points) & Stufen nach GER (CEFR Level) \\
                \hline
                75--100         & B1                           \\
                35--74          & A2                           \\
                0--34           & unter A2 (below A2)          \\
                \hline
            \end{tabular}
        \end{table}
    % \bibliographystyle{plain}
    % \bibliography{/home/donkarlo/Dropbox/repo/shared_research/bibliography/bibliography.bib}
\end{document}